\appendixsection{Модель данных и хранилища}
\label{app:data-model}

\subsection{Назначение приложения}

Приложение фиксирует целевую модель хранения данных платформы после перехода
template-registry на MongoDB. Полные DDL- и миграционные артефакты оставлены
в исходном коде репозитория.

\subsection{Гибридная модель хранения}

В текущей версии используется три хранилища:
\begin{enumerate}
    \item PostgreSQL для операционных контуров событий, запусков доставки и истории;
    \item MongoDB для шаблонов уведомлений и их версий;
    \item Redis для коммуникационного профиля получателя.
\end{enumerate}

Распределение контуров по хранилищам приведено в таблице~\ref{tab:appendix-storage-layout}.

\begin{table}[h]
    \centering
    \footnotesize
    \renewcommand{\arraystretch}{1.15}
    \caption{Распределение данных по хранилищам}
    \label{tab:appendix-storage-layout}
    \begin{tabular}{|L{0.22\textwidth}|L{0.32\textwidth}|L{0.36\textwidth}|}
        \hline
        \textbf{Хранилище} & \textbf{Контуры} & \textbf{Назначение} \\
        \hline
        PostgreSQL & nf\_fac, nf\_sched, nf\_mail, nf\_sms, nf\_hist & Операционные данные, outbox/inbox-журналы, статусы и история доставки. \\
        \hline
        MongoDB & template-registry (`templates`) & Шаблоны, версии шаблонов, канальные части шаблона, активная версия. \\
        \hline
        Redis & profile-consent & Профиль получателя, ограничения каналов, адресные атрибуты. \\
        \hline
    \end{tabular}
\end{table}

\subsection{Коллекция шаблонов в MongoDB}

Сервис template-registry использует одну коллекцию `templates`.
Каждый документ включает:
\begin{enumerate}
    \item метаданные шаблона (`templateId`, `clientId`, `name`, `description`, `status`);
    \item номер активной версии `activeVersion`;
    \item idempotency-ключ создания `createIdempotencyKey`;
    \item массив `versions`, содержащий версии и канальные представления;
    \item технические поля (`createdAt`, `updatedAt`, `schemaVersion`, `revision`).
\end{enumerate}

\subsection{Индексы и ограничения MongoDB}

Для коллекции `templates` применяются индексы:
\begin{enumerate}
    \item уникальный составной индекс `(clientId, templateId)`;
    \item уникальный sparse-индекс `(clientId, createIdempotencyKey)`;
    \item индекс `(clientId, updatedAt)` для `ListTemplates`.
\end{enumerate}

Целостность состава версии (допустимые каналы, непрерывность версий,
корректность `activeVersion`) обеспечивается на уровне прикладной логики
сервиса template-registry.

\subsection{Пример документа шаблона}

Ниже приведён обезличенный пример документа коллекции `templates`,
используемого в тестовом контуре:

\begin{verbatim}
{
  "templateId": "9fbf6f6a-7ef1-4f5a-bf9f-95cb9d0d4e2f",
  "clientId": "2e08a6f8-1fc5-4eb7-bb9a-c54c5a3028af",
  "name": "order-status",
  "description": "Шаблон уведомления о статусе заказа",
  "status": "TEMPLATE_STATUS_PUBLISHED",
  "activeVersion": 3,
  "createIdempotencyKey": "create-order-status-v1",
  "schemaVersion": 1,
  "versions": [
    {
      "version": 3,
      "engine": "TEMPLATE_ENGINE_HANDLEBARS",
      "contents": [
        {
          "channel": "CHANNEL_EMAIL",
          "subject": "Заказ {{orderId}} обновлён",
          "body": "Здравствуйте, {{name}}. Новый статус: {{status}}."
        },
        {
          "channel": "CHANNEL_SMS",
          "subject": "",
          "body": "Заказ {{orderId}}: статус {{status}}."
        }
      ],
      "createdAt": "2026-04-13T11:10:34Z"
    }
  ],
  "createdAt": "2026-04-13T09:00:11Z",
  "updatedAt": "2026-04-13T11:10:34Z"
}
\end{verbatim}

\subsection{Ключевые ограничения данных}

Критичные ограничения модели:
\begin{enumerate}
    \item уникальность события по паре \texttt{client\_id} и \texttt{idempotency\_key} в командном контуре;
    \item уникальность запуска доставки в границах события;
    \item уникальность факта канальной доставки для ключевой бизнес-комбинации;
    \item уникальность шаблона по паре `(clientId, templateId)` в MongoDB;
    \item уникальность идемпотентного создания шаблона по `(clientId, createIdempotencyKey)`.
\end{enumerate}

\subsection{Служебные журналы}

Для контролируемого обмена и обработки используются два класса служебных данных:
\begin{enumerate}
    \item исходящие журналы интеграционных событий;
    \item входящие журналы потреблённых сообщений.
\end{enumerate}

Их совместное применение обеспечивает трассируемость и восстановление обработки
после частичных сбоев без потери состояния.

\subsection{Redis-модель профиля получателя}

Профиль получателя хранится в Redis-хеше и включает:
\begin{enumerate}
    \item признак активности;
    \item предпочитаемый канал;
    \item атрибуты доступности и адресации для email, SMS и push;
    \item метку актуальности профиля.
\end{enumerate}

Доступ к модели профиля выполняется только через выделенный сервисный контракт,
что исключает прямую зависимость остальных сервисов от внутреннего формата
хранилища.
